\documentclass[12pt,nofootinbib,aps,amsmath,amssymb,preprint,showpacs]{revtex4}

\usepackage[T1]{fontenc} \usepackage[latin1]{inputenc}
\usepackage{amsmath,color,ulem} \usepackage{graphicx}
\usepackage{epsf,bm} \usepackage{amssymb}
\newcommand\underrel[2]{\mathrel{\mathop{#2}\limits_{#1}}}
\usepackage{tabulary}
\newcolumntype{K}[1]{>{\centering\arraybackslash}p{#1}}
\newcommand\underrela[2]{\mathrel{\mathop{#2}\limits_{#1}}}
\usepackage{multirow}
\begin{document}

\title{Procedural maze generation algorithm based on mechanics from Blue Prince}
\author{Dennis van den Bogaard$^\dagger$}
\author{Tigo Schippers$^\dagger$}
 
\affiliation{
 $^\dagger$Department of Information and Computing Sciences, Utrecht University, Princetonplein 5, 3584 CC Utrecht, The Netherlands\\
}

\date{\today} 

\begin{abstract} 
\end{abstract}

\pacs{05.10.Gg, 05.10.Ln, 05.40.-a, 05.50.+q, 05.70.Jk}
\maketitle
 
\section{Introduction \label{sce1}}

\section{Methods\label{sec2}}

\section{Results \label{sec3}}

\section{Discussion \label{sec4}}

\section{Conclusion \label{sec5}}

\begin{thebibliography}{11}

\bibitem {} H. E. Stanley, {\it Phase transitions and critical
  phenomena}, Clarendon Press, Oxford (1971).
\bibitem {} Gabrovšek, P. (2019). Analysis of Maze Generation Algorithms. {\it IPSI Trans. Internet Res}. 6. https://ipsitransactions.org/journals/papers/tir/2019jan/p5.pdf 

\end{thebibliography}
\end{document}
